\section{Иерархия вложенных кристаллов и соответствие с физикой твёрдого тела}

В классической физике твёрдого тела атомы образуют пространственную решётку Браве с периодом~$A$,
элементарная область вокруг узла задаётся многогранником Вигнера--Зейтца,
а коллективные возбуждения описываются фононами с квазиимпульсом в первой зоне Бриллюэна~$|k|\le \pi/A$.
В решёточной жидкости микроуровня~$\Mlayer$ эта же кристаллографическая структура лежит в самом основании пространства-времени,
порождая иерархию вложенных кристаллов (\cref{ch:matter,ch:si-sm}):
\begin{enumerate}[leftmargin=*]
\item \textbf{Первый (планковский) кристалл ($\Mlayer$, шаг $a = \caStepSpace = \ell_P$):}
  узлы гранецентрированной кубической решётки несут элементарные ромбические додекаэдры Вороного объёма $\mathrm{d}V = \caKappaGeom\,a^3$,
  заполненные кипящей спинорной жидкостью $z\in\mathbb{Z}_{N_{\mathrm{ring}}}[i]^2$ с бинарной плотностью занятости ядра $b_{\Mlayer}\in\{0,1\}$.
  Первая зона Бриллюэна с границей $k_{\max} = \pi/a$ и частотой Дебая $\omega_D = 2\pi/\htick$ исключает существование волн короче~$a$.
\item \textbf{Второй (атомный) кристалл ($\Tlayer$, шаг $A = K\cdot a$, $K\in\mathbb{N}$):}
  устойчивые связанные состояния барионных треугольников~$K_3$ и электронных оболочек образуют макро-узлы с занятостью $b_{\mathrm{atom}}\in\{0,1\}$,
  атомные ячейки Вигнера--Зейтца которых имеют ту же форму ромбического додекаэдра ГЦК и вмещают целое число $N_{\mathrm{pack}}(K) = V_{\mathrm{WS}}(K\cdot a)/\mathrm{d}V = K^3\in\mathbb{Z}$ планковских ячеек.
\end{enumerate}

Благодаря самоподобию ромбического додекаэдра ГЦК и сокращению масштаба~$K$ в кванте действия $s_{\mathrm{floor}}(K) = (p_0/K)(K\cdot a) = \ladderAction$,
законы симметрии каналов кубооктаэдра $(V,E,\square,\triangle)$ переносятся с планковского кристалла на атомный и молекулярный кристаллы без введения новых параметров.
